\chapter{如何建構虛構世界 --- Kate Messner}

Messner 一開始便做對了一件方法論上的事。她不是從「世界建構」的定義講起，而是從那些已經運作良好的世界講起。Gandalf 可以死去又歸來；在 Matrix 裡，技能可以被下載；Cheshire Cat 可以拋接自己的頭；一場 Quidditch 比賽，直到 Golden Snitch 被抓住才會結束；而終極問題的答案是 \(42\)。只有在這一長串被記住的規則展示過後，她才告訴我們：把它們連在一起的不是「什麼都可以發生」，而是「被允許發生的一切，都受穩定規則支配」。接著，講座便沿著一條乾淨的序列展開：例子、原則、與日常現實的對照、讀者對書中世界的掌握、更深一層的疑問，即頁面上的痕跡如何變成活生生的經驗，最後則是作者如何實際建構一個世界，並讓故事從其中浮現。

\section{先看例子：有規則的世界}

講座的開場是累積性的。每一個例子都在提醒我們：陌生之物之所以能夠令人信服，不是因為它奇特，而是因為它穩定。Gandalf 的身體是凡俗的，而其靈體是不朽的。Matrix 有一套讓能力加速下載的機制。Cheshire Cat 的邏輯荒謬，卻不散亂。Quidditch 有明確的終止條件。甚至 Douglas Adams 的 \(42\) 之所以有力，也正因為它聽起來像是從某個世界的規律運作中被計算出來的答案，而那個世界的問答方式並不服從我們自己的規則。

若用一種簡潔的筆記形式記下這段開場，便是
\begin{align}
\text{凡俗身體} + \text{不朽靈體} &\Rightarrow \text{Gandalf 可以死去又回來},\\
\text{接入} + \text{破解程式碼} &\Rightarrow \text{在數秒內學會直升機飛行},\\
\text{Snitch 被抓住} &\iff \text{Quidditch 比賽結束},\\
\text{終極答案} &= 42.
\end{align}
這些並不是嚴格意義上的方程式，而是局部規則陳述。Messner 的重點是：讀者正是透過這類規則來記住一個世界的。驚奇之所以可理解，是因為它有形式。

這就使講座得以提出它的第一個明確抽象：
\begin{equation}
\text{一致的規則} \Rightarrow \text{可信性} + \text{可理解性} + \text{可探索性}.
\end{equation}
一個只是讓人吃驚的世界是薄的；一個世界之所以變得可探索，是因為某一次事件教會我們如何閱讀下一次事件。因此，講座從例子開始，不是為了在真正論證前先娛樂我們，而是為了讓真正的論證變得可見。

\subsection{Question \& Answer}

\paragraph{Question.} 什麼讓一個被發明出來的世界顯得連貫，而不是隨機？

\paragraph{Answer.} 是它的奇觀受到約束。一旦讀者可以推知這個世界允許什麼、禁止什麼、又會反覆出現什麼，它便不再是一堆零散的新奇，而成了一個有結構的領域。

\section{真實法則、頁面法則，以及作者真正建造的是什麼}

接著，Messner 做出一個更尖銳的比較。在我們的世界裡，重力讓 Harry Potter 的書安穩地待在書架上；這一點我們知道是真的。但在小說頁面之中，一旦巫師世界被建立起來，重力便可以被修改、被暫停，或者被另一條規則所回應。決定性的區分，不是「有法則」與「無法則」，而是「一個有法則的領域」與「另一個有法則的領域」。

對於 Messner 所說的 ``a spectrum of physical and societal rules''，一種有用的壓縮寫法如下。

\begin{definition}
我們可以把一個虛構世界示意地寫成
\begin{equation}
W = (R_{\mathrm{phys}},\,R_{\mathrm{soc}}),
\end{equation}
其中 \(W\) 表示這個世界，\(R_{\mathrm{phys}}\) 表示其物理規則，\(R_{\mathrm{soc}}\) 表示其社會規則。這是筆記撰寫者加入的記號，而不是講座中的記號，但它忠實捕捉了講座實際上的結構。
\end{definition}

於是，那個關於重力的對照便可以審慎地寫成
\begin{equation}
g_{\text{現實世界}} \neq g_{\text{巫師頁面世界}}.
\end{equation}
這樣的速記不是要假裝 Messner 在講物理，而是要把她散文中已然在做的工作說得更精確：區分讀者所處世界的法則，與書中世界的法則。

這也就是為什麼她接著能立刻把尺度放大。科幻與奇幻的作者並不是只選擇一塊背景；他們建造的是一整套可理解性的基礎設施：
\begin{equation}
\text{規則} \to \text{地圖} \to \text{譜系} \to \text{語言} \to \text{文化} \to \text{宇宙}.
\end{equation}
這些箭頭標記的是被清晰表述的結構不斷增加。一條被改動的法則只給我們一個點子；一套彼此相連的法則、歷史與制度，才給我們一個世界。到了這裡，講座便準備好進一步追問：世界不只是什麼，它還會對讀者做什麼。

\section{當讀者理解了書中的世界}

Messner 接下來的主張，是整場講座第一個真正的回報。當世界建構做得足夠好時，讀者對虛構世界及其規則的理解，可以與角色一樣深，甚至比讀者對書外世界的理解還更深。這不是一句附帶的旁白，而是製造出整場講座核心張力的命題。因為如果這件事真會發生，那我們就必須追問：它究竟是怎麼發生的？

若用一種簡潔的筆記式比較，我們可以寫成
\begin{equation}
\text{讀者對書中世界的掌握} \gtrsim \text{讀者對外部世界的掌握}.
\end{equation}
符號 \(\gtrsim\) 只是對 Messner 那句 ``just as well or even better'' 的速記，而不是講座中真正說出的記號，但它讓這個主張的力度保持在眼前。

講座裡那張經過驗證的圖像，在這裡尤其有用。它不是方程式，而是一種空間對照。前景裡是讀者和攤開的書，一個小小的故事世界從書頁中浮起，外部世界則在後景裡作為另一塊社會場域存在。這張圖示出了那句話的意思：書中的世界，可以暫時成為更可讀的那個領域。

\begin{figure}[h]
\centering
\includegraphics[width=0.82\textwidth]{lecture_09_figure_03.png}
\caption{讀者、書中世界，以及書外世界。}
\end{figure}

由於這張截圖的修辭性大於符號性，因此在旁邊放上一張清理過的示意圖會有幫助。我們保留原圖作為證據，而把 TikZ 重繪作為詮釋。

\begin{figure}[h]
\centering
\begin{tikzpicture}[x=1cm,y=1cm]
\node[draw, rounded corners, minimum width=2.0cm, minimum height=0.9cm] (reader) at (0,0) {\small 讀者};
\node[draw, minimum width=2.2cm, minimum height=0.9cm] (book) at (2.9,0) {\small 打開的書};
\node[draw, rounded corners, minimum width=2.1cm, minimum height=0.9cm] (inside) at (5.9,0.95) {\small 書中世界};
\node[draw, rounded corners, minimum width=2.5cm, minimum height=0.9cm] (outside) at (6.2,-1.35) {\small 外部世界};
\draw[->, thick] (reader) -- (book);
\draw[->, thick] (book) -- (inside);
\draw[dashed, thick] (book.east) -- (outside.west);
\end{tikzpicture}
\caption{對講座中「書中世界」與「書外世界」對照的示意性重建。}
\end{figure}

真正重要的，不在於重繪圖形的精確幾何，而在於講座的那個直覺：閱讀可以生產出一個其規則比日常生活更向我們敞開的領域。

\subsection{Question \& Answer}

\paragraph{Question.} 一本書裡的世界怎麼可能比現實更可讀？

\paragraph{Answer.} 因為書中世界是為了可理解性而被建構的。它的規則是經過挑選、重複並彼此支撐的。現實更豐富，但也更嘈雜。因此，講座的意思是：一個建構良好的虛構世界，對讀者而言反而可能是更可處理的系統。

\section{可是，怎麼做到的？彎曲線條怎麼變成世界？}

正是在這裡，Messner 暫停了整個向前的推進，提出講座中最大的問題：可是，怎麼做到的？頁面上的彎曲線條，怎麼能反射光進入眼睛，向大腦發送訊號，被解碼為敘事，進而在情緒上打動我們，最後甚至改變我們在讀完最後一筆之後回到世界時的觀看方式？

講座本身的鏈條可以簡潔地寫成
\begin{equation}
\text{頁面上的彎曲線條} \to \text{發往大腦的訊號} \to \text{被解碼的敘事} \to \text{情緒／思想／行動} \to \text{改變了的視角}.
\end{equation}
這裡，我們應當保留講座本身的節奏。Messner 並沒有把問題壓縮成單純的資訊傳遞而已。她明白地把這條鏈，從知覺擴大到感受，從感受擴大到行動，再從行動擴大到視角的變化。閱讀的終點不只是理解，而是轉變，無論那轉變有多細微。

\paragraph{Worked derivation.} 若把它拆成步驟，講座中的機制便是
\begin{enumerate}
\item 一些痕跡被放到頁面上。
\item 頁面反射出的光抵達眼睛。
\item 神經系統把這些輸入轉成訊號。
\item 心智把這些有模式的訊號解碼為敘事。
\item 敘事召喚出情緒、思想與想像。
\item 讀者進入虛構世界，把它當成一個被活過的領域。
\item 在離開這本書之後，讀者對日常世界的看法也可能有所不同。
\end{enumerate}

這是一條功能性的推演，而不是一套完整的閱讀科學。Messner 對此十分謹慎。她明白說過，自己並不確定是否有人真正完全知道答案。這個讓步很重要。講座並沒有假裝要先把這個謎解決掉，才能繼續談技藝；它讓這個謎保留下來，然後仍然轉向實作。

\subsection{Question \& Answer}

\paragraph{Question.} 頁面上的痕跡，到底怎麼變成讀者心中一個被活過的世界？

\paragraph{Answer.} 講座不是用定理，而是用一條鏈來回答。痕跡變成訊號，訊號變成被解碼出的秩序，被解碼的秩序變成敘事，而敘事又足以強烈地召喚感受，進而產生沉浸與回返。它的精確機制仍有部分晦暗；但可操作的序列已經很清楚了。

\section{從謎走向方法}

一旦講座容許這個謎保持晦暗，它便改變了尺度。Messner 說，奇幻世界每天都在被創造：在我們心裡、在電腦裡，甚至在餐巾紙上。這裡的實際前提是：想像力，加上願意住在想像力裡的意志，就足以讓事情開始。接著，講座便下沉到作者的方法之中。

第一步是歷史性的：
\begin{equation}
\text{過去事件} \to \text{時間軸} \to \text{世界當下的結構}.
\end{equation}
這個主張比表面上更強。我們不是從裝飾開始，而是從追問開始：究竟發生過什麼？發生的順序是什麼？長遠後果又是什麼？因此，一條時間軸並不是附加材料，而是對「這個世界為何被安排成現在這樣」的第一個形式答案。

從這裡開始，Messner 透過一組組有順序的問題往前推。首先，在宏觀尺度上，我們問的是法律、權力、信念與價值：
\begin{equation}
(\text{規則},\ \text{政府},\ \text{權力},\ \text{信念},\ \text{價值}).
\end{equation}
然後，在日常生活的尺度上，我們下沉到被活著的質地之中：
\begin{equation}
(\text{天氣},\ \text{家／工作／學校},\ \text{食物},\ \text{玩樂},\ \text{年輕／年老},\ \text{植物／動物},\ \text{科技},\ \text{交通},\ \text{溝通},\ \text{資訊}).
\end{equation}
這樣的排列在教學上極其準確。我們先從構成性法則走到社會結構，然後才走到日常生活。一個世界若只給我們口號，還不能住人；它是在我們知道人們吃什麼、如何移動、害怕什麼、教導什麼、懲罰什麼，以及早晨迎接他們的是什麼樣的天氣時，才真正變得可居住。

\paragraph{Worked example.} Messner 的方法可以被改寫成一條有紀律的世界描述更新規則。
\begin{enumerate}
\item 指定這個世界的過去，並把它組織成一條時間軸。
\item 決定有哪些物理規則與社會規則正在生效。
\item 追問由誰統治、誰握有權力，以及人們信什麼。
\item 再追問這個社會珍視什麼，以及如何處理違反者。
\item 往日常現實下探：天氣、工作、學校、食物、遊戲、年齡、生態與科技。
\item 一直做下去，直到這個世界不再只是清單，而成為一個地方。
\end{enumerate}
講座並不是要我們做出鉅細靡遺的行政表格，而是要求一種足夠有結構的注意力，直到這個世界開始能夠連貫地回應我們。

\section{從世界走向故事}

直到最後，Messner 才揭露講座真正的目的地。世界建構並不是終點；故事才是。一旦我們認識這個世界的程度，已經接近我們希望讀者能認識它的程度，我們便把角色放進去，讓他們自由活動，然後追問接下來會發生什麼。這時，世界就不再只是背景，而成了因果力量的一部分。

Messner 的最後壓縮可以寫成
\begin{equation}
\text{世界} + \text{角色} \Rightarrow \text{被形塑的個體} \Rightarrow \text{衝突} \Rightarrow \text{故事}.
\end{equation}
這條方程式值得立在結尾，因為它重新解釋了前面所有內容。那些規則從來不只是裝飾，那條時間軸也從來不只是檔案，問卷式的盤點更從來不只是描述而已。它們全都只是為了壓力做準備。世界決定了它的居民可以渴望什麼、害怕什麼、違反什麼、成為什麼，也因此決定了哪些類型的衝突可能發生。

我們可以把這個點說得更直白些。沒有人的世界，只會給我們一份資料夾；沒有世界的人，只會給我們抽象物。故事開始於：被結構化的人物，與被結構化的世界相遇，並在其限制之下被迫行動。

\subsection{Question \& Answer}

\paragraph{Question.} 什麼時候世界建構不再只是背景，而開始成為故事？

\paragraph{Answer.} 當世界的結構開始支配選擇時，它就成為故事。一旦角色必須活在這個世界的規則、價值、科技、懲罰與允許之下，衝突便出現了。到那時，世界不再是風景，而是戲劇中一個積極運作的變項。

\section{Summary}

Messner 的講座為虛構世界提供了一套緊湊而真實的數學。強而有力的創造世界，必須受規則支配。它們的法則可以與日常生活截然不同，但這種差異本身也必須連貫。當建構成功時，讀者會以驚人的清晰度學會一本書內部的世界。而這件事反過來又引出更深的謎：頁面上的痕跡，究竟如何變成被活過的經驗。講座容許這個謎仍舊部分不可解，卻也在技藝層次上給出實際答案：建構世界的過去，指定它的法則與制度，下沉到日常生活，然後讓角色進入那套已成形的結構。就在那一刻，世界建構便不再只是準備，而開始變成故事。